\documentclass{article}
\usepackage[utf8]{inputenc}
\usepackage{geometry}
\usepackage{hyperref}
\usepackage{enumitem}
\usepackage{xcolor}

\geometry{a4paper, margin=1in}

\title{Phase 2: Video Structure \& Script Outline}
\author{Kaushik Vada}
\date{}

\begin{document}

\maketitle

\section*{1. The Hook (0:00 - 0:45)}
\textbf{Visual:} You on camera.

\textbf{Action:} Hold up a modern CPU (or picture of one) and a stick of RAM.

\textbf{Script Concept:} "This processor runs at 4 Gigahertz—that's 4 billion cycles a second. This RAM is huge, but compared to the CPU, it’s a tortoise. If the CPU had to wait for RAM for every single instruction, your supercomputer would feel like a calculator from 1980. So, how do we fake speed? We cheat using something called Cache."

\textbf{Diagram:} [Insert Diagram Here]

\section*{2. The "Why" \& The Problem (0:45 - 1:30)}
\textbf{Visual:} Animation of a fast car (CPU) stuck in traffic behind a tractor (RAM).

\textbf{Explanation:} Explain the "Memory Wall". The CPU is too fast for the memory. We need a middleman.

\section*{3. The Solution: Multi-Level Hierarchy (1:30 - 4:00)}
\textbf{Visual:} Split screen. Left side: Technical diagram of CPU Cores, L1, L2, L3. Right side: The Kitchen Animation.

\textbf{Explanation:}
\begin{itemize}
    \item \textbf{L1 (Level 1):} "This is the Chef's cutting board. It's built inside the core. It holds the variables you are using right now. It's tiny (kilobytes) but instant."
    \item \textbf{L2 (Level 2):} "The Countertop. If the Chef (CPU) needs garlic and it's not on the board (L1 Miss), they check here. It's bigger (megabytes) but slightly slower."
    \item \textbf{L3 (Level 3):} "The Walk-in Pantry. This is special. It is usually shared across multiple cores. If Core 1 and Core 2 both need the same data, they can find it here without going to RAM."
\end{itemize}

\section*{4. How They Work Together: The "Miss" Chain (4:00 - 6:00)}
\textbf{Visual:} A flowchart animation showing a "Data Request."

\textbf{Scenario:} The CPU needs the "Health Bar" value for a video game character.

\begin{itemize}
    \item \textbf{Check L1:} Is it there? NO (Cache Miss). \textit{Visual: Chef looks at empty cutting board.}
    \item \textbf{Check L2:} Is it there? NO (Cache Miss). \textit{Visual: Chef checks the fridge. Empty.}
    \item \textbf{Check L3:} Is it there? YES (Cache Hit). \textit{Visual: Chef runs to the pantry, finds the ingredient, brings it to the fridge (L2), then to the board (L1).}
\end{itemize}

\textbf{The Key Takeaway:} "Data is copied up the chain. The CPU only ever directly touches L1. L2 and L3 are just waiting rooms to keep L1 full."

\section*{5. Real World Application (6:00 - 7:30)}
\textbf{Visual:} Screen recording of a heavy application or game loading.

\textbf{Application:} Explain why "Loading Screens" exist (filling the RAM) vs. why games stutter (waiting for data to move from RAM to Cache).

\textbf{Analogy connection:} "A stutter in your game is the Chef waiting for the delivery truck (RAM) because the Pantry (L3) was empty."

\section*{6. Conclusion \& Self-Intro (7:30 - End)}
\textbf{Recap:} Fast memory is expensive, so we use small layers of it.

\textbf{Self-Intro:} "Hi, I'm Kaushik Vada, a student in CS161. I hope this explains why your computer is faster than it has any right to be."

\end{document}
